\section{Treewidth}

% lax begin Lax261.Treewidth
\begin{definition}[Treewidth]
  A \emph{tree decomposition} of a graph $G = (V, E)$ is a pair
  $(T, \{X_t\}_{t \in V(T)})$ where $T$ is a tree and each $X_t \subseteq V$
  is a \emph{bag}, such that
  \begin{itemize}
    \item every vertex of $G$ lies in some bag,
    \item for every edge $uv \in E$ some bag contains both $u$ and $v$, and
    \item for every $v \in V$ the bags containing $v$ induce a subtree of $T$.
  \end{itemize}
  The \emph{width} of the decomposition is $\max_t (|X_t| - 1)$, and the
  \emph{treewidth} $\tw(G)$ is the minimum width over all tree
  decompositions of $G$.
\end{definition}
% lax end

The subtree condition is the one that makes dynamic programming over a
decomposition work: it guarantees that the bags containing a fixed vertex
form a connected region of the tree.

% lax begin Lax261.Treewidth
\begin{theorem}
  A graph $G$ satisfies $\tw(G) \le 1$ if and only if $G$ is a forest.
\end{theorem}

% lax begin Lax261Proofs.Q
\begin{proof}
  If $G$ is a forest, take $T = G$ with singleton and edge bags; every bag
  has size at most two, so the width is at most one.  Conversely, a cycle
  forces two non-adjacent vertices into one bag, a contradiction.
\end{proof}
% lax end

% lax end
